(1) 代入 \(\phi=Ae^{i(kx-\omega t)}\)，对流时间微分给 \(-i(\omega-Uk)\)，空间微分给 \(ik\)，于是 \((\omega-Uk)^2=c^2k^2\)。有两支 \(\omega_\pm=(U\pm c)k\)，流体共动参考系中频率为 \(\omega'=\omega-Uk=\pm ck\)。频率平移是 Galilean 对流效应，而不是声速本身被改变。

(2) 取各支符号与 \(k\) 约定一致，实验室相速度 \(v_{p,\pm}=\omega_\pm/k=U\pm c\)，因其色散线性，群速度 \(v_{g,\pm}=d\omega_\pm/dk=U\pm c\)。亚声速 \(|U|<c\) 下有一支朝下游，一支朝上游传播；若顺流 \(U>c\)，上游支的实验室速度也朝下游，说明上游声信号不能穿过超声速对流区域。

(3) 方程可因式分解为 \([\partial_t+(U+c)\partial_x][\partial_t+(U-c)\partial_x]\phi=0\)，由于两个算子系数常数且可交换，一般解为 \(\phi=f(x-(U+c)t)+g(x-(U-c)t)\)。若给定初始位移及其时间导数，分别对 \(f,g\) 施加和差约束即可确定传播分量；两条特征线斜率分别为 \(U+c,U-c\)。

(4) 若实验室声源固定角频率 \(\omega_s>0\)，下游波数 \(k_+=\omega_s/(U+c)\)，上游取实验室传播速度 \(U-c<0\)，其正向坐标波数 \(k_-=\omega_s/(U-c)<0\)。当 \(U\to c^-\)，上游波长 \(2\pi/|k_-|\to0\)，表明在理想无色散线性模型中逆流传播越来越难；实际短波近似、粘性与流动非均匀性终将限制该极限。
